% Full title: Chinese Translation of Corrected Representation Identity
\section{修正后的表示恒等式}
\label{app:corrected-identity}

这里记录互补场论证中的波数分配，以避免一种常见的不匹配。若 $w^{\mathrm e}$ 满足背景方程 $(\Delta+k^2)w^{\mathrm e}=0$，且 $v^{\mathrm i}$ 是互补内部 $k$ 方程的解，则准周期消去数据对为
\begin{equation}
  -S^{\mathrm{qp}}_{k,\beta}
  [\gamma_N^{\mathrm i}v^{\mathrm i}]
  +D^{\mathrm{qp}}_{k,\beta}
  [\gamma_D^{\mathrm i}v^{\mathrm i}].
  \label{eq:corrected-exterior-pair}
\end{equation}
\eqref{eq:corrected-exterior-pair} 中的两个算子都使用背景波数 $k$。夹杂表示使用自由空间数据对
\begin{equation}
  -S_{k^0}[\gamma_N^{\mathrm i}w^{\mathrm i}]
  +D_{k^0}[\gamma_D^{\mathrm i}w^{\mathrm i}],
  \qquad k^0=n^0k.
  \label{eq:corrected-interior-pair}
\end{equation}
若把 $D^{\mathrm{qp}}_{k^0,\beta}$ 混入 \eqref{eq:corrected-exterior-pair}，则外部 $k$-Helmholtz 消去恒等式失效。\eqref{eq:corrected-exterior-pair} 与 \eqref{eq:corrected-interior-pair} 符合 \eqref{eq:center-reconstruction} 中的波数分离，也符合 \cite[Appendix~A]{BarnettGreengard2010} 的结构。
